\chapter{Introduction}
\label{ch:introduction}

\section{Background}
\label{sec:introduction-background}

Vehicular networks are becoming an important part of connected mobility, intelligent transportation systems, and future wireless infrastructure. In vehicle-to-infrastructure (V2I) communication, roadside base stations (BSs) are expected to provide high-rate and reliable links for vehicles whose positions, velocities, and blockage conditions vary rapidly over time. The use of millimetre-wave and other high-frequency bands can provide large communication bandwidth, but it also makes the communication link more dependent on narrow beams, accurate beam alignment, and timely beam maintenance \cite{xu2022computervisionaidedmmwave,jiang2022lidaraidedfuturebeam,liu2020radarassistedpredictivebeamformingvehicular}.

In conventional high-frequency cellular systems, beam management is commonly supported by beam sweeping, synchronization signal block (SSB) measurements, channel state information reference signal (CSI-RS) measurements, and subsequent channel state information (CSI) feedback. These mechanisms allow the BS and the user equipment to identify suitable transmit and receive beams and to support later precoding decisions. However, when the users are vehicles moving through dynamic road environments, the validity of a beam measurement can be short-lived. A beam that is suitable at one frame may become suboptimal after a short displacement, a lane change, or the emergence of a blockage. Performing frequent full beam sweeps can reduce this uncertainty, but it consumes time-frequency resources that would otherwise be used for data transmission \cite{alrabeiah2019millimeterwavebasestations,charan2021visionpositionmultimodalbeamprediction,charan2023camerabasedmmwavebeam}.

Roadside infrastructure also has a natural sensing role. A BS can be equipped with cameras and integrated sensing and communication (ISAC) capability, while nearby sensing nodes (SNs) can provide additional views of the road using distributed cameras. These sensing observations are not merely auxiliary context: they can reveal vehicle locations, motion patterns, occlusion relationships, and road geometry before or alongside communication measurements. In this thesis, the sensing infrastructure is treated as a source of spatial information that can assist beam management in a BS-centric vehicular network. The service BS collects its own multiview camera observations, its BS-side ISAC point cloud, and the multiview camera observations from nearby SNs. These inputs are then fused in a bird's-eye-view (BEV) representation to support vehicle detection and per-target beam prediction.

The central motivation is that beam management in vehicular networks should not be considered only as a communication measurement problem. If the network can identify which detected vehicle corresponds to which communication user, and if it can maintain this identity over time, sensing information can be used to narrow the beam search space and reduce the frequency of full beam sweeping. This motivates a joint workflow that couples multimodal sensing, target-user association, tracking, and system-level communication evaluation \cite{alrabeiah2019millimeterwavebasestations,demirhan2021radaraided6gbeam,jiang2022lidaraidedfuturebeam,liu2022learningbasedpredictivebeamformingintegrated}.

\section{Challenges}
\label{sec:introduction-challenges}

The first challenge is the fast temporal variation of the wireless link. In a road environment, a vehicle may change its relative angle and distance to the serving BS over a small number of frames. The received beam power distribution over a discrete BS codebook can therefore change quickly, particularly in high-frequency bands where narrow beams are used. A sensing-assisted system must predict useful candidate beams from visual and ISAC observations while accounting for the fact that the communication channel and the sensing scene are both time-varying.

The second challenge is blockage and measurement uncertainty. Vehicles, roadside objects, and surrounding infrastructure can occlude line-of-sight or dominant reflection paths. At the same time, sensing observations can be incomplete because of camera field-of-view limits, synchronization errors, calibration errors, or sparse ISAC point clouds. Communication measurements can also be noisy or outdated. A practical beam management workflow must therefore tolerate uncertain observations across both sensing and communication modalities rather than assuming that either side provides a perfect state estimate.

The third challenge is multi-vehicle identity ambiguity. A detected sensing target is not automatically a communication user identity. In a scene with multiple vehicles, the sensing network may output several target detections and several per-target beam distributions, while the communication system observes beam power spectra associated with communication users. The system must determine which detected target corresponds to which user before the predicted beams can be used for user-specific communication decisions. This target-user association problem is essential in multi-vehicle sensing-assisted beam management, but it is distinct from the beam index prediction problem itself \cite{charan2023camerabasedmmwavebeam,imran2024environmentsemanticcommunicationenabling,zhou2020trackingobjectspoints}.

The fourth challenge is maintaining association over time. Repeating a full SSB beam sweep at every frame can re-establish user-specific beam observations, but it increases overhead. Conversely, relying only on a previous association can lead to beam mismatch when the vehicle moves, when tracking drifts, or when the scene becomes occluded. This motivates a tracking mechanism that predicts the state of associated users, matches new BEV detections to existing tracks, and decides when additional communication measurement is needed. In this thesis, the tracking component is based on an interacting multiple model (IMM) filter that combines multiple vehicle-motion hypotheses for temporally continuous identity maintenance.

These challenges indicate that a sensing-assisted vehicular beam management system should be evaluated beyond Top-K beam accuracy alone. Incorrect association, tracking failure, beam mismatch, and beam training overhead can all affect the final communication performance. For this reason, this thesis connects sensing outputs to a simplified system-level communication evaluation that includes beam training overhead, zero-forcing (ZF) precoding, effective signal-to-interference-plus-noise ratio (SINR), spectral efficiency, and overhead-adjusted throughput.

\section{Research Objective}
\label{sec:introduction-objective}

The objective of this thesis is to develop and evaluate a BS-centric multimodal distributed sensing-assisted framework for low-overhead and trackable beam management in vehicular networks. The framework uses the service BS as the coordination point. For each decision frame, the BS collects its local multiview camera observations, its BS-side ISAC point cloud, and the multiview camera observations from a selected set of nearby SNs. These multimodal observations are fused in a BEV representation to detect vehicles and to predict a per-target beam power distribution over the BS beam codebook.

The thesis focuses on four connected tasks. The first task is BEV-based vehicle detection from multimodal distributed sensing observations. The second task is per-target beam prediction, where each detected vehicle is assigned a predicted beam power distribution and a set of Top-K candidate beams. The third task is target-user association, where the initial SSB beam sweep power spectrum of each communication user is matched with the predicted beam distribution of detected sensing targets. The fourth task is tracking-assisted beam maintenance, where associated users are maintained over subsequent frames by combining BEV detections with an IMM tracking state.

The communication objective is not to reproduce a complete 5G New Radio protocol stack. Instead, this thesis adopts a simplified and interpretable system-level abstraction. In this abstraction, sensing-assisted Top-K beam prediction can reduce beam training overhead when the correct beam is included in the candidate set, while an incorrect candidate set introduces a beam mismatch penalty that affects subsequent CSI or precoding decisions. The resulting effective SINR, spectral efficiency, and overhead-adjusted throughput are then used to study how sensing accuracy, association reliability, tracking continuity, and measurement overhead interact at the communication-system level.

The resulting study is methodological and system-oriented: it defines the sensing-to-communication workflow, specifies the association and tracking interface, and evaluates how beam prediction, identity maintenance, and communication overhead interact in the proposed vehicular setting.

\section{Contributions}
\label{sec:introduction-contributions}

This thesis is organized around three contributions.

\begin{enumerate}
    \item A BS-centric multimodal distributed sensing-assisted beam management framework is developed for vehicular networks. The framework uses the serving BS as the coordination point and fuses BS-side multiview camera observations, BS-side ISAC point clouds, and camera observations from nearby distributed SNs. The fused BEV representation supports vehicle target detection and per-target beam power distribution prediction, from which Top-K candidate beams are obtained for later beam management decisions.

    \item A target-user association and tracking workflow is designed for multi-vehicle sensing-assisted communication. During initial access, the workflow uses the SSB beam sweep power spectrum from the communication side and the predicted per-target beam distribution from the sensing side to associate detected sensing targets with communication user identities. In later frames, an IMM tracking module and BEV detection association are used to maintain user identity and support temporally continuous beam selection while reducing reliance on repeated full beam sweeping.

    \item A multimodal vehicular simulation system and system-level evaluation method are constructed to connect sensing outputs with communication metrics. The simulation workflow links road traffic generation, multisensor observation generation, scene reconstruction, ray-tracing-based channel generation, beam power labels, and simplified communication evaluation. The evaluation considers Top-K beam prediction, target-user association, tracking error, beam training overhead, beam mismatch penalty, ZF precoding SINR, spectral efficiency, and overhead-adjusted throughput.
\end{enumerate}

The contributions are stated at the framework, association-and-tracking, and system-evaluation levels; quantitative performance claims are made only where they are supported by the experimental evaluation.

\section{Thesis Organization}
\label{sec:introduction-organization}

The remainder of this thesis is organized as follows. Chapter~\ref{ch:background} reviews the background and related work on vehicular beam management, sensing-assisted beam prediction, ISAC predictive beamforming, BEV multimodal fusion, cooperative perception, target-user association, and tracking. It also identifies the research gap addressed by this thesis.

Chapter~\ref{ch:system-model} presents the system model and problem formulation. It defines the roadside vehicular network, BSs, distributed SNs, vehicle users, multimodal sensing observations, target-user association variables, tracking states, and the communication and beam management model used for system-level evaluation.

Chapter~\ref{ch:methodology} describes the proposed multimodal sensing-assisted beam management framework. It introduces the BS-centric data flow, BEV-based multimodal fusion, vehicle detection head, per-target beam prediction head, initial target-user association procedure, IMM tracking module, and the interface between sensing outputs and communication evaluation.

Chapter~\ref{ch:experiments} describes the multimodal simulation system and experimental evaluation. It covers scenario and dataset construction, baselines, evaluation metrics, main results, ablation studies, and discussion of limitations.

Finally, Chapter~\ref{ch:conclusion} concludes the thesis and discusses future work. It summarizes the findings supported by the thesis, clarifies the main limitations of the simulation and communication abstraction, and outlines possible extensions toward more realistic protocol-level simulation, robust multimodal sensing, and broader cooperative vehicular network settings.
